\subsection{其他选择准则}\label{sec:entropy-selection-criteria}

\PDEIndexMark{idx:other-selection-criteria-introduction}Lax 激波条件并非从守恒律方程组的可能弱解中选出"物理上合理"解的唯一可行选择准则. 本节介绍一些竞争性的条件, 并描述它们之间的若干关系.

\subsubsection{熵条件}\label{subsec:entropy-condition}

首先介绍的另一种选择准则称为熵条件. 它源自热力学第二定律, 在这里将该定律推广到还包括机械系统与热系统之外的物理系统. 这一条件的关键是存在一个由 \eqref{eq:conservation-law-system} 导出的附加守恒律.

\begin{Definition}
\label{def:entropy-flux-pair}
\textbf{熵/熵通量对}\footnote{尽管这一术语与热力学中任何合理的定义都相矛盾, 它却已相当标准. 也许更好的术语是自由能/能量通量, 但经过这么多年后再试图改变似乎已毫无希望. 因此, 学习守恒律的学生在研究这一主题时, 只得接受忘掉所学过的全部热力学知识这一事实. 对这一术语负有责任的始作俑者尚未站出来承认.}是满足
\begin{equation}\label{eq:entropy-flux-pair-definition}
\nabla F=\nabla U\cdot\nabla\bm{f}.
\end{equation}
的函数对 $(U,F):\mathbb{R}^n\to\mathbb{R}^2$.
\end{Definition}

由定义与链式法则立即可知, 若 $\bm{u}$ 为 \eqref{eq:conservation-law-system} 的经典解, 则
\begin{equation}\label{eq:classical-entropy-conservation}
U(\bm{u})_t+F(\bm{u})_x=0.
\end{equation}
当然, 如 Remark~\ref{rem:weak-solution-variable-change} 所指出的, 方程 \eqref{eq:conservation-law-system} 的弱解不一定满足 \eqref{eq:classical-entropy-conservation}.

\begin{Definition}
\label{def:entropy-condition}
若存在一个熵/熵通量对, 其中 $\bm{u}\mapsto U(\bm{u})$ 为凸函数, 使得对每个在开上半平面中具有紧支集的\emph{非负} $C^1$ 检验函数 $\phi$ 都有
\begin{equation}\label{eq:entropy-weak-inequality}
-\int\!\!\int(U(\bm{u})\phi_t+F(\bm{u})\phi_x)\,\mathrm{d}x\,\mathrm{d}t\leq0
\end{equation}
则称 \eqref{eq:conservation-law-system} 的弱解满足\textbf{熵条件}.
\end{Definition}

\begin{Remark}
\label{rem:entropy-distribution-interpretation}
满足 \eqref{eq:entropy-weak-inequality} 的熵/熵通量对, 常被描述为在\emph{分布意义下}满足不等式
\begin{equation}\label{eq:entropy-distribution-inequality}
U(\bm{u})_t+F(\bm{u})_x\leq0
\end{equation}
注意, 若一个熵/熵通量对在经典意义下满足 \eqref{eq:entropy-distribution-inequality}, 则可将该不等式乘以非负检验函数, 再利用平面上的 Green 定理导出 \eqref{eq:entropy-weak-inequality}.
\end{Remark}

\begin{Remark}
\label{rem:entropy-free-energy}
若 $\bm{u}$ 满足 \eqref{eq:entropy-distribution-inequality} 并具有紧支集, 则将 \eqref{eq:entropy-distribution-inequality} 对 $x$ 积分, 得到
\begin{equation}\label{eq:total-entropy-time-derivative}
\int_{-\infty}^\infty U_t\,\mathrm{d}x\leq0.
\end{equation}
因此, 量
\[
\int_{-\infty}^\infty U\,\mathrm{d}x
\]
(再次强调, 最好将其理解为自由能而非熵) 是时间的递减函数. 能量只能损失这一思想, 可视为热力学第二定律的某个版本.
\end{Remark}

以下引理将熵条件转化为分片连续弱解的跳跃条件.

\begin{Lemma}
\label{lem:entropy-jump-condition}
假设 $\bm{u}$ 为满足熵条件的 \eqref{eq:conservation-law-system} 的分片连续弱解. 假设 $\bm{u}$ 沿一条斜率为 $s$ 的激波曲线发生跳跃间断. 那么
\begin{equation}\label{eq:entropy-jump-condition}
s[U(\bm{u})]-[F(\bm{u})]\geq0.
\end{equation}

\end{Lemma}

证明留给读者 (参见 Problem~\ref{ex:entropy-jump-condition-proof}). 可通过修改 Rankine-Hugoniot 条件的证明得到这一证明.

以下定理指出, 对真正非线性的方程组, 对于弱激波, 熵条件与 Lax 激波条件等价.

\begin{Theorem}
\label{thm:entropy-lax-small-shock-equivalence}
假设守恒律方程组 \eqref{eq:conservation-law-system} 严格双曲且真正非线性. 此外, 假设存在熵/熵通量对 $(U,F)$, 其中 $U$ 强凸. 设 $\bm{u}$ 为 \eqref{eq:conservation-law-system} 的弱解, 含有一个以速度 $s$ 传播且两侧值为 $\bm{u}^l$ 和 $\bm{u}^r$ 的激波, 并且 $\bm{u}^l-\bm{u}^r$ 充分小. 那么 $\bm{u}$ 在该激波处满足 Lax 激波条件, 当且仅当熵跳跃条件 \eqref{eq:entropy-jump-condition} 以严格不等式成立.
\end{Theorem}

\begin{Proof}
由于 $\bm{u}$ 是弱解, $\bm{u}^l$, $\bm{u}^r$ 与 $s$ 必须满足 Rankine-Hugoniot 条件. 因此, 由 Theorem~\ref{thm:shock-curves-existence}, $\bm{u}^r$ 必须位于通过 $\bm{u}^l$ 的 $n$ 条激波曲线之一上; 即对某个 $k=1,2,\ldots,n$ 和某个 $\epsilon\in I_k$, 有 $\bm{u}^r=\hat{\bm{u}}_k(\epsilon)$ 和 $s=\hat s_k(\epsilon)$. 再由 Theorem~\ref{thm:shock-curves-existence}, Lax 激波条件满足当且仅当 $\epsilon<0$. 此外, 熵跳跃条件成立当且仅当
\begin{equation}\label{eq:entropy-shock-curve-functional}
E(\epsilon):=\hat s_k(\epsilon)[U(\bm{u}^l)-U(\hat{\bm{u}}_k(\epsilon))]-[F(\bm{u}^l)-F(\hat{\bm{u}}_k(\epsilon))]<0.
\end{equation}
现在用"撇号"表示对 $\epsilon$ 求导, 并证明 $E(0)=E'(0)=E''(0)=0$ 及 $E'''(0)>0$. 这说明, 当 $\epsilon$ 充分小时, $E$ 为负当且仅当 $\epsilon$ 为负, 从而完成证明. 利用 $\bm{u}^l=\hat{\bm{u}}_k(0)$, 得到
\begin{equation}\label{eq:entropy-shock-functional-at-zero}
E(0)=\hat s_k(0)[U(\bm{u}^l)-U(\hat{\bm{u}}_k(0))]-[F(\bm{u}^l)-F(\hat{\bm{u}}_k(0))]=0.
\end{equation}
本可用 $\hat s_k(0)=\lambda_k$ 和 $\bm{u}_k'(0)=\bm{r}_k$ 直接计算 $E'(0)$, 但改为对 Rankine-Hugoniot 条件求导, 得到
\begin{equation}\label{eq:shock-rh-parameter-derivative}
\hat s_k'(\epsilon)[\bm{u}^l-\hat{\bm{u}}_k(\epsilon)]-\hat s_k(\epsilon)\hat{\bm{u}}_k'(\epsilon)=-\bm{f}(\hat{\bm{u}}_k(\epsilon))'.
\end{equation}
再利用 $\nabla F=\nabla U\cdot\nabla\bm{f}$, 得到
\begin{equation}\label{eq:entropy-shock-functional-first-derivative}
E'(\epsilon)=\hat s_k'(\epsilon)\{[U(\bm{u}^l)-U(\hat{\bm{u}}_k(\epsilon))]-\nabla U(\hat{\bm{u}}_k(\epsilon))\cdot[\bm{u}^l-\hat{\bm{u}}_k(\epsilon)]\},
\end{equation}
由此容易看出 $E'(0)=0$. 直接对其求导即得 $E''(0)=0$. 计算 $E'''(0)$ 时, 许多项与前面计算中的项以相同方式变为零, 但留下一个有趣的项:
\begin{equation}\label{eq:entropy-shock-functional-third-derivative}
E'''(0)=\hat s_k'(0)[\hat{\bm{u}}_k'(0)^T\nabla^2U(\hat{\bm{u}}_k(0))\hat{\bm{u}}_k'(0)]
\end{equation}
其中 $\nabla^2U$ 为 $U$ 的二阶梯度或 Hessian 矩阵. 现在, 由 Theorem~\ref{thm:shock-curves-existence} 有 $\hat s_k'(0)=1/2$, 并且由于 $U$ 严格凸, 它的 Hessian 矩阵正定. 因此 $E'''>0$, 定理得证.
\end{Proof}
